\documentclass[a4paper,11pt]{ltjsarticle}
\usepackage{luatexja}
\usepackage{amsmath,amssymb,amsthm}
\usepackage{longtable,booktabs,array,calc}
\usepackage{hyperref}
\usepackage[margin=25mm]{geometry}

\newtheorem{theorem}{定理}[section]
\newtheorem{proposition}[theorem]{命題}
\newtheorem{lemma}[theorem]{補題}
\newtheorem{corollary}[theorem]{系}
\newtheorem{definition}[theorem]{定義}
\newtheorem{remark}[theorem]{注意}
\newtheorem{observation}[theorem]{観察}

\providecommand{\tightlist}{\setlength{\itemsep}{0pt}\setlength{\parskip}{0pt}}
\newcounter{none}

\begin{document}

\section{6次元超直方体の集合構造とその組合せ論的性質}\label{ux6b21ux5143ux8d85ux76f4ux65b9ux4f53ux306eux96c6ux5408ux69cbux9020ux3068ux305dux306eux7d44ux5408ux305bux8ad6ux7684ux6027ux8cea}

\subsection{------ スピン分類・集合の数え上げ・符号付き面積
------}\label{ux30b9ux30d4ux30f3ux5206ux985eux96c6ux5408ux306eux6570ux3048ux4e0aux3052ux7b26ux53f7ux4ed8ux304dux9762ux7a4d}

\textbf{著者}: 木原 範昭 (WF System Co., Ltd.) \textbf{日付}: 2026年4月
\textbf{DOI}: 10.5281/zenodo.19646652

\begin{center}\rule{0.5\linewidth}{0.5pt}\end{center}

\subsection{0. 本稿の目的}\label{ux672cux7a3fux306eux76eeux7684}

本稿は、5 次元超直方体に離散ラベル軸 \(Q\) を追加した 6
次元構造を定義し、その集合構造から導かれる組合せ論的性質を述べる。

本体（§1--§9）では、本モデルの数学的構造と、そこから純粋に組合せ論的に導出される帰結のみを主張する。標準模型との対応は
§10
で「解釈の一例」として提示するが、これは本論文の主張ではない。モデルの集合構造が標準模型の粒子分類とどのように対応付けうるかを示す一例に過ぎない。

\textbf{本モデルの外部仮定について}: 本稿は 5
次元超直方体を出発点とする。なぜ 5 次元であるか、なぜ 5 軸のうち 3
軸を後に「空間軸」として扱うかは、\textbf{本モデルの外部仮定}であり、モデル内部では正当化されない。観測される時空が
3 次元空間 + 1
次元時間であるという経験的事実への対応として、この次元設定を採用している。

\begin{center}\rule{0.5\linewidth}{0.5pt}\end{center}

\subsection{1. 6
次元構造の定義}\label{ux6b21ux5143ux69cbux9020ux306eux5b9aux7fa9}

\subsubsection{1.1 定義}\label{ux5b9aux7fa9}

本稿で扱う 6 次元構造は以下で定義される。

\begin{itemize}
\tightlist
\item
  \textbf{幾何学的構造}: 5
  次元超直方体の組合せ構造（頂点、辺、面、3-cell、4-facet、5-cell 本体）
\item
  \textbf{ラベル軸}: 4 値の離散ラベル軸 \(Q \in \{0, r, g, b\}\)
\end{itemize}

5 次元超直方体の頂点座標は \((x_1, x_2, x_3, x_4, x_5)\)
で表記し、各座標は\textbf{任意の実数値をとる}。特定の命題が集合構造において座標値に条件を課す場合は、その命題の中で局所的に宣言する。

第 6 軸 \(Q\)
は幾何学的な座標ではなく、各状態に付随する離散ラベル値である。値域は
\(\{0, r, g, b\}\) の 4 値とする。\(Q\)
軸の追加によって超直方体の組合せ構造（頂点数、辺数、面数など）は変化しない。

ただし、スピンを決定する双ベクトルの分類においては、\(Q\) は他の 5
軸と同様に 1
本の軸として向きの数え上げに参加する。すなわち、\(e_t \wedge e_Q\) や
\(e_R \wedge e_Q\) のような組合せは、幾何学的な 2-face
ではなく、軸の組合せとして定義される抽象的な双ベクトル対である。

\subsubsection{\texorpdfstring{1.2 \(k\)
次元超直方体の要素数}{1.2 k 次元超直方体の要素数}}\label{k-ux6b21ux5143ux8d85ux76f4ux65b9ux4f53ux306eux8981ux7d20ux6570}

\(k\) 次元超直方体の \(j\) 次元面（\(j\)-face）の個数を \(f_j(k)\)
とする。各要素数を低次元から書き下すことで、一般式を読者が検算できる形で提示する。

\textbf{頂点数 \(f_0(k)\)}: 各軸について両端の 2
状態を区別する組合せ数として

{\def\LTcaptype{none} % do not increment counter
\begin{longtable}[]{@{}ccccccc@{}}
\toprule\noalign{}
\(k\) & 0 & 1 & 2 & 3 & 4 & 5 \\
\midrule\noalign{}
\endhead
\bottomrule\noalign{}
\endlastfoot
\(f_0(k)\) & 1 & 2 & 4 & 8 & 16 & 32 \\
\end{longtable}
}

\[f_0(k) = 2^k\]

\textbf{辺の数 \(f_1(k)\)}: 1 本の軸方向に伸び、残り \((k-1)\)
本の軸の両端位置で決まる

{\def\LTcaptype{none} % do not increment counter
\begin{longtable}[]{@{}ccccccc@{}}
\toprule\noalign{}
\(k\) & 0 & 1 & 2 & 3 & 4 & 5 \\
\midrule\noalign{}
\endhead
\bottomrule\noalign{}
\endlastfoot
\(f_1(k)\) & 0 & 1 & 4 & 12 & 32 & 80 \\
\end{longtable}
}

\[f_1(k) = k \cdot 2^{k-1}\]

検算: \(k=3\) で \(3 \cdot 4 = 12\)（立方体の辺 12 本）、\(k=5\) で
\(5 \cdot 16 = 80\)。

\textbf{正方形面の数 \(f_2(k)\)}: 2 本の軸で張られる面

{\def\LTcaptype{none} % do not increment counter
\begin{longtable}[]{@{}ccccccc@{}}
\toprule\noalign{}
\(k\) & 0 & 1 & 2 & 3 & 4 & 5 \\
\midrule\noalign{}
\endhead
\bottomrule\noalign{}
\endlastfoot
\(f_2(k)\) & 0 & 0 & 1 & 6 & 24 & 80 \\
\end{longtable}
}

\[f_2(k) = \binom{k}{2} \cdot 2^{k-2}\]

検算: \(k=3\) で \(3 \cdot 2 = 6\)（立方体の面 6 枚）、\(k=5\) で
\(10 \cdot 8 = 80\)。

\textbf{立方体胞（3-cell）の数 \(f_3(k)\)}: 3 本の軸で張られる立方体胞

{\def\LTcaptype{none} % do not increment counter
\begin{longtable}[]{@{}ccccccc@{}}
\toprule\noalign{}
\(k\) & 0 & 1 & 2 & 3 & 4 & 5 \\
\midrule\noalign{}
\endhead
\bottomrule\noalign{}
\endlastfoot
\(f_3(k)\) & 0 & 0 & 0 & 1 & 8 & 40 \\
\end{longtable}
}

\[f_3(k) = \binom{k}{3} \cdot 2^{k-3}\]

検算: \(k=4\) で \(4 \cdot 2 = 8\)、\(k=5\) で \(10 \cdot 4 = 40\)。

\textbf{4-facet の数 \(f_4(k)\)}: 4 本の軸で張られる 4 次元セル

{\def\LTcaptype{none} % do not increment counter
\begin{longtable}[]{@{}ccccccc@{}}
\toprule\noalign{}
\(k\) & 0 & 1 & 2 & 3 & 4 & 5 \\
\midrule\noalign{}
\endhead
\bottomrule\noalign{}
\endlastfoot
\(f_4(k)\) & 0 & 0 & 0 & 0 & 1 & 10 \\
\end{longtable}
}

\[f_4(k) = \binom{k}{4} \cdot 2^{k-4}\]

検算: \(k=5\) で \(5 \cdot 2 = 10\)。

\textbf{5-cell 本体の数 \(f_5(k)\)}: \(k=5\) の場合

\[f_5(5) = \binom{5}{5} \cdot 2^0 = 1\]

\subsubsection{一般式}\label{ux4e00ux822cux5f0f}

以上をまとめると、\(k\) 次元超直方体の \(j\) 次元面の個数は

\[f_j(k) = \binom{k}{j} \cdot 2^{k-j} \qquad (0 \leq j \leq k)\]

で与えられる。これは、\(k\) 本の軸から \(j\) 本を選んで \(j\)-face
を張り（\(\binom{k}{j}\) 通り）、残り \((k-j)\)
本の軸について各両端のどちらに位置するかを指定する（\(2^{k-j}\)
通り）ことに対応する。

\textbf{検算（オイラー標数）}:
\(\sum_{j=0}^{k} (-1)^j f_j(k) = (2-1)^k = 1\)
により、超直方体が可縮空間であることと整合する。

\subsubsection{\texorpdfstring{本稿で用いる値（\(k=5\)）}{本稿で用いる値（k=5）}}\label{ux672cux7a3fux3067ux7528ux3044ux308bux5024k5}

{\def\LTcaptype{none} % do not increment counter
\begin{longtable}[]{@{}clc@{}}
\toprule\noalign{}
\(j\) & 要素 & \(f_j(5)\) \\
\midrule\noalign{}
\endhead
\bottomrule\noalign{}
\endlastfoot
0 & 頂点 & 32 \\
1 & 辺 & 80 \\
2 & 正方形面 & 80 \\
3 & 立方体胞（3-cell） & 40 \\
4 & 4-facet & 10 \\
5 & 5-cell 本体 & 1 \\
\end{longtable}
}

\subsubsection{1.3
軸の対称性と集合符号の定義}\label{ux8ef8ux306eux5bfeux79f0ux6027ux3068ux96c6ux5408ux7b26ux53f7ux306eux5b9aux7fa9}

\paragraph{1.3.1 5
軸の完全対称性}\label{ux8ef8ux306eux5b8cux5168ux5bfeux79f0ux6027}

\textbf{命題 1.1（5 軸の完全対称性）}. 5 次元超直方体の 5
本の幾何学的軸は、組合せ構造上すべて対等である。任意の軸の置換は超直方体の自己同型写像を与え、§1.2
の要素数 \(f_j(k)\) および後続の幾何学的命題に関して 5
軸は区別されない。

\textbf{証明}. \(k\) 次元超直方体の組合せ構造の自己同型群は超八面体群
\(B_k\) であり、任意の軸置換を含む。したがって 5
軸は組合せ的に完全対称である。\(\square\)

\paragraph{1.3.2 軸の分類}\label{ux8ef8ux306eux5206ux985e}

\textbf{定義 1.2（軸の表記ラベル）}. 5 本の軸
\(x_1, x_2, x_3, x_4, x_5\) に対して以下の表記ラベルを与える。

\begin{itemize}
\tightlist
\item
  \(x_1, x_2, x_3\) を \(x, y, z\) と表記する
\item
  \(x_4\) を \(t\) と表記する
\item
  \(x_5\) を \(R\) と表記する
\end{itemize}

命題 1.1 により 5
軸は組合せ構造上対等であるから、どの軸にどのラベルを付すかは定義の選択に過ぎない。これらのラベルは集合符号の定義（§1.3.3）およびスピンの定義（§2）の表記の前提として用いる。

\textbf{\(x, y, z, t, R\)
はあくまで表記上のラベルであり、幾何学的軸としての区別を意味しない。}
物理的解釈（空間・時間・曲率など）を各軸に割り当てる操作は
§10（解釈の一例）で行うが、これはモデル外部からの解釈であり、§1
の幾何学的構造とは独立である。

\paragraph{1.3.3
集合符号の厳密な定義}\label{ux96c6ux5408ux7b26ux53f7ux306eux53b3ux5bc6ux306aux5b9aux7fa9}

各幾何学的軸の成分
\(x_i\)（\(i = 1, \ldots, 5\)）は、その軸方向の座標値が満たすべき\textbf{範囲制約}を表す記号である。本稿では以下の
5 種類の記号を用いる。

{\def\LTcaptype{none} % do not increment counter
\begin{longtable}[]{@{}cll@{}}
\toprule\noalign{}
記号 & 意味 & 座標値の範囲 \\
\midrule\noalign{}
\endhead
\bottomrule\noalign{}
\endlastfoot
\(+\) & 正値 & \(x_i > 0\) \\
\(-\) & 負値 & \(x_i < 0\) \\
\(0\) & 原点 & \(x_i = 0\) \\
\(\pm\) & 非零値 & \(x_i \neq 0\)（正または負） \\
\(*\) & 任意 & \(x_i\) は任意の実数 \\
\end{longtable}
}

\(Q\) 軸についても同様に、値域を表す記号を用いる。

{\def\LTcaptype{none} % do not increment counter
\begin{longtable}[]{@{}
  >{\centering\arraybackslash}p{(\linewidth - 4\tabcolsep) * \real{0.2778}}
  >{\raggedright\arraybackslash}p{(\linewidth - 4\tabcolsep) * \real{0.3333}}
  >{\raggedright\arraybackslash}p{(\linewidth - 4\tabcolsep) * \real{0.3889}}@{}}
\toprule\noalign{}
\begin{minipage}[b]{\linewidth}\centering
記号
\end{minipage} & \begin{minipage}[b]{\linewidth}\raggedright
意味
\end{minipage} & \begin{minipage}[b]{\linewidth}\raggedright
\(Q\) の値
\end{minipage} \\
\midrule\noalign{}
\endhead
\bottomrule\noalign{}
\endlastfoot
\(0\) & 色中性 & \(Q = 0\) \\
\(r, g, b\) & 指定色 & \(Q\) は指定された色ラベル \\
\(r/g/b\) & 3 色のいずれか & \(Q \in \{r, g, b\}\) \\
\(\bar{Q}\) & 反色 &
\(r \leftrightarrow \bar{r}, g \leftrightarrow \bar{g}, b \leftrightarrow \bar{b}, 0 \leftrightarrow 0\) \\
\(*\) & 任意 & \(Q\) は \(\{0, r, g, b\}\) のいずれか \\
\end{longtable}
}

\textbf{注記}.
集合符号は座標値そのものではなく、座標値が満たすべき\textbf{条件}を指定する記号である。§1.1
で述べた通り、本モデルの座標値は任意の実数値をとり得るが、特定の命題や集合構造においては、該当する軸の座標値に上記の範囲制約を課す。

\textbf{例}. 集合 \((+, +, 0, 0, 0, 0)\)
は、\(x_1 > 0, x_2 > 0, x_3 = 0, x_4 = 0, x_5 = 0, Q = 0\)
を満たす状態の集合を指す。

\begin{center}\rule{0.5\linewidth}{0.5pt}\end{center}

\subsection{2. スピンの定義}\label{ux30b9ux30d4ux30f3ux306eux5b9aux7fa9}

\textbf{定義 2.1（スピン）}. 集合 \((x, y, z, t, R, Q)\)
に対して、\(t, R, Q\) 軸のうち集合符号が \(0\) でない軸の本数を \(n\)
とする。スピン \(s\) を以下で定義する。

\begin{itemize}
\tightlist
\item
  \(xyz\) 軸のうち 2 本が非零で 1 本が \(0\) のとき:
  \(s = n + 1/2\)（フェルミオン）
\item
  それ以外のとき: \(s = n\)（ボソン）
\end{itemize}

\(t, R, Q\) は 3 軸であるから、\(n\) のとりうる値は \(0, 1, 2, 3\) の 4
通りである。したがって、本モデルにおけるスピン \(s\)
のとりうる値は以下の 8 種となる。

{\def\LTcaptype{none} % do not increment counter
\begin{longtable}[]{@{}ccc@{}}
\toprule\noalign{}
\(n\) & ボソン \(s = n\) & フェルミオン \(s = n + 1/2\) \\
\midrule\noalign{}
\endhead
\bottomrule\noalign{}
\endlastfoot
0 & \(s = 0\) & \(s = 1/2\) \\
1 & \(s = 1\) & \(s = 3/2\) \\
2 & \(s = 2\) & \(s = 5/2\) \\
3 & \(s = 3\) & \(s = 7/2\) \\
\end{longtable}
}

\textbf{注記}.
本定義は本モデル内部の組合せ論的な定義である。ここで「ボソン」「フェルミオン」と呼ぶのは本モデル内の用語であり、物理的なスピン角運動量、統計性、\((-1)^{2s}\)
変換性などとの対応は §10（解釈の一例）で扱う。

\begin{center}\rule{0.5\linewidth}{0.5pt}\end{center}

\subsection{\texorpdfstring{3. \(Q\)
ラベル間遷移の数え上げ}{3. Q ラベル間遷移の数え上げ}}\label{q-ux30e9ux30d9ux30ebux9593ux9077ux79fbux306eux6570ux3048ux4e0aux3052}

\(Q\) 軸の 3 つの非零値 \(\{r, g, b\}\) の間の遷移を考える。遷移は
\((Q_\text{初期}, Q_\text{終期})\) の対で記述され、3 値から 3
値への遷移の総数は

\[3 \times 3 = 9\]

通りである。

このうち、対角成分の 1 つの線形結合（\(r\bar{r} + g\bar{g} + b\bar{b}\)
に比例する状態）は、すべての色ラベルに対して同じ作用を持ち、ラベル遷移を媒介しない。この
1 状態を除いた

\[3 \times 3 - 1 = 8\]

を独立な遷移状態の数とする。物理的命名は §10 で行う。

\begin{center}\rule{0.5\linewidth}{0.5pt}\end{center}

\subsection{4. スピン 1
集合の数え上げ}\label{ux30b9ux30d4ux30f3-1-ux96c6ux5408ux306eux6570ux3048ux4e0aux3052}

スピン 1 の集合は定義 2.1 により \(n = 1\)、すなわち \(t, R, Q\)
軸のうち集合符号が \(0\) でない軸が 1 本であり、\(xyz\) 軸はすべて \(0\)
の状態である。

非零の軸が \(t, R, Q\)
のいずれかに応じて、集合の状態数は以下のように分類される。

{\def\LTcaptype{none} % do not increment counter
\begin{longtable}[]{@{}
  >{\raggedright\arraybackslash}p{(\linewidth - 4\tabcolsep) * \real{0.3200}}
  >{\raggedright\arraybackslash}p{(\linewidth - 4\tabcolsep) * \real{0.3600}}
  >{\centering\arraybackslash}p{(\linewidth - 4\tabcolsep) * \real{0.3200}}@{}}
\toprule\noalign{}
\begin{minipage}[b]{\linewidth}\raggedright
非零の軸
\end{minipage} & \begin{minipage}[b]{\linewidth}\raggedright
集合形式
\end{minipage} & \begin{minipage}[b]{\linewidth}\centering
状態数
\end{minipage} \\
\midrule\noalign{}
\endhead
\bottomrule\noalign{}
\endlastfoot
\(t\) & \((0, 0, 0, \pm, 0, 0)\) & 2（符号 \(\pm\) の 2 通り） \\
\(R\) & \((0, 0, 0, 0, \pm, 0)\) & 2（符号 \(\pm\) の 2 通り） \\
\(Q\) & \((0, 0, 0, 0, 0, Q_\text{初期} \to Q_\text{終期})\) & 9（§3
による） \\
\end{longtable}
}

\textbf{合計}: \(2 + 2 + 9 = 13\) 状態（うち 1 状態を除く解釈については
§10 参照）

\begin{center}\rule{0.5\linewidth}{0.5pt}\end{center}

\subsection{5. スピン 2
集合の数え上げ}\label{ux30b9ux30d4ux30f3-2-ux96c6ux5408ux306eux6570ux3048ux4e0aux3052}

スピン 2 の集合は定義 2.1 により \(n = 2\)、すなわち \(t, R, Q\)
軸のうち集合符号が \(0\) でない軸が 2 本であり、\(xyz\) 軸はすべて \(0\)
の状態である。

\(t, R, Q\) の 3 軸から 2 軸を選ぶ組合せは \(\binom{3}{2} = 3\)
通りであり、非零となる 2 軸の組合せに応じて以下のように分類される。

{\def\LTcaptype{none} % do not increment counter
\begin{longtable}[]{@{}ll@{}}
\toprule\noalign{}
非零の 2 軸 & 集合形式 \\
\midrule\noalign{}
\endhead
\bottomrule\noalign{}
\endlastfoot
\(t, R\) & \((0, 0, 0, \pm, \pm, 0)\) \\
\(t, Q\) & \((0, 0, 0, \pm, 0, Q)\) \\
\(R, Q\) & \((0, 0, 0, 0, \pm, Q)\) \\
\end{longtable}
}

したがって、本モデルにおけるスピン 2 集合の型は \textbf{3 通り} である。

\begin{center}\rule{0.5\linewidth}{0.5pt}\end{center}

\subsection{6. スピン 3
集合の数え上げ}\label{ux30b9ux30d4ux30f3-3-ux96c6ux5408ux306eux6570ux3048ux4e0aux3052}

スピン 3 の集合は定義 2.1 により \(n = 3\)、すなわち \(t, R, Q\)
軸のすべてが非零であり、\(xyz\) 軸はすべて \(0\) の状態である。

\(t, R, Q\) の 3 軸すべてが非零となる組合せは \(\binom{3}{3} = 1\)
通りである。

{\def\LTcaptype{none} % do not increment counter
\begin{longtable}[]{@{}ll@{}}
\toprule\noalign{}
非零の 3 軸 & 集合形式 \\
\midrule\noalign{}
\endhead
\bottomrule\noalign{}
\endlastfoot
\(t, R, Q\) & \((0, 0, 0, \pm, \pm, Q)\)（\(Q \neq 0\)） \\
\end{longtable}
}

したがって、本モデルにおけるスピン 3 集合の型は \textbf{1 通り} である。

\begin{center}\rule{0.5\linewidth}{0.5pt}\end{center}

\subsection{7.
非零軸の符号積}\label{ux975eux96f6ux8ef8ux306eux7b26ux53f7ux7a4d}

\textbf{定義 7.1（符号積）}. ボソン集合において、\(t, R, Q\)
軸のうち集合符号が \(0\) でない \(n\) 本の軸がそれぞれ \(+\) または
\(-\) の符号をとるとき、これらの符号の積を\textbf{符号積} \(P(n)\)
と定義する。

\[P(n) = (-1)^n\]

{\def\LTcaptype{none} % do not increment counter
\begin{longtable}[]{@{}cc@{}}
\toprule\noalign{}
\(n\) & \(P(n)\) \\
\midrule\noalign{}
\endhead
\bottomrule\noalign{}
\endlastfoot
0 & \(+1\) \\
1 & \(-1\) \\
2 & \(+1\) \\
3 & \(-1\) \\
\end{longtable}
}

\textbf{命題 7.1}. \(P(n)\) は \(n\) の偶奇のみに依存する。\(n\)
が偶数のとき \(P(n) = +1\)、\(n\) が奇数のとき \(P(n) = -1\) である。

\textbf{注記}. \(n = 0\)（符号反転なし）と \(n = 2\)（二重反転）は同じ
\(P = +1\) を与えるが、その代数的構造は異なる。\(n = 0\)
は恒等変換であり、\(n = 2\) は符号反転の自己逆元である。

\(P(n)\) の物理的解釈（力の方向性との対応）は §10（解釈の一例）で扱う。

\begin{center}\rule{0.5\linewidth}{0.5pt}\end{center}

\subsection{\texorpdfstring{8. \(R\)
軸を含む正方形面の符号付き面積}{8. R 軸を含む正方形面の符号付き面積}}\label{r-ux8ef8ux3092ux542bux3080ux6b63ux65b9ux5f62ux9762ux306eux7b26ux53f7ux4ed8ux304dux9762ux7a4d}

\subsubsection{8.1
符号付き面積の定義}\label{ux7b26ux53f7ux4ed8ux304dux9762ux7a4dux306eux5b9aux7fa9}

\textbf{定義 8.1（符号付き面積）}. 集合構造 \(\sigma\)
の\textbf{符号付き面積} \(M(\sigma)\) を、\(R\)
軸を含む全正方形面の集合符号と面積の積の総和として定義する。

\[M(\sigma) = \sum_{R \text{ 軸を含む正方形面 } f} \mathrm{sgn}(f) \cdot \mathrm{Area}(f)\]

\(M(\sigma)\)
は本モデルの基本的な保存量である。物理的な解釈（質量との類似性など）は
§10 で扱う。

\subsubsection{\texorpdfstring{8.2 \(R\)
軸を含む面の分類}{8.2 R 軸を含む面の分類}}\label{r-ux8ef8ux3092ux542bux3080ux9762ux306eux5206ux985e}

80
個の正方形面のうち、\(R\)（\(= x_5\)）軸を含む面は以下のように分類される。

{\def\LTcaptype{none} % do not increment counter
\begin{longtable}[]{@{}llcc@{}}
\toprule\noalign{}
面の型 & 双ベクトル & 型数 & 面数 \\
\midrule\noalign{}
\endhead
\bottomrule\noalign{}
\endlastfoot
空間\(-R\) & \(e_i \wedge e_5\) & 3 & \(3 \times 8 = 24\) \\
\(t-R\) & \(e_4 \wedge e_5\) & 1 & 8 \\
\textbf{合計} & & 4 & \textbf{32} \\
\end{longtable}
}

\(R\) 軸を含む面は全 80 面中 32 面であり、全体の \(2/5\) を占める。

\subsubsection{\texorpdfstring{8.3 \(M(\sigma) = 0\)
となる集合の存在}{8.3 M(\textbackslash sigma) = 0 となる集合の存在}}\label{msigma-0-ux3068ux306aux308bux96c6ux5408ux306eux5b58ux5728}

\textbf{命題 8.1}. 集合構造 \(\sigma\)
の中には、\(M(\sigma) = 0\)（符号付き面積が完全に相殺する）となるものが存在する。

\(M(\sigma) = 0\) となる集合と物理的な対応の解釈は §10 で扱う。

\begin{center}\rule{0.5\linewidth}{0.5pt}\end{center}

\subsection{9. 本体の結論}\label{ux672cux4f53ux306eux7d50ux8ad6}

本稿では、5 次元超直方体に離散ラベル軸 \(Q\) を追加した 6
次元構造を定義し、その集合構造から以下の組合せ論的性質を導出した。

\begin{enumerate}
\def\labelenumi{\arabic{enumi}.}
\tightlist
\item
  \textbf{超直方体の要素数}（§1.2）: \(k\) 次元超直方体の \(j\)
  次元面の個数 \(f_j(k) = \binom{k}{j} \cdot 2^{k-j}\)。\(k = 5\)
  の場合、32 頂点、80 辺、80 正方形面、40 立方体胞、10 超直方体。
\item
  \textbf{5 軸の完全対称性}（命題 1.1）: 5
  本の幾何学的軸は組合せ構造上すべて対等であり、任意の軸置換は超直方体の自己同型写像を与える。
\item
  \textbf{軸の表記ラベル}（定義 1.2）: 5 軸に \(x, y, z, t, R\)
  の表記ラベルを与える。
\item
  \textbf{集合符号の厳密な定義}（§1.3.3）: 各軸に対する範囲制約を
  \(+, -, 0, \pm, *\) の 5 記号で表現する。
\item
  \textbf{スピンの定義}（定義 2.1）: 集合 \((x, y, z, t, R, Q)\)
  からスピン \(s\) を、\(t, R, Q\) 軸の非零本数
  \(n\)（\(= 0, 1, 2, 3\)）と \(xyz\) 軸の構造に基づいて 8
  種（\(s = 0, 1/2, 1, 3/2, 2, 5/2, 3, 7/2\)）に分類する。
\item
  \textbf{\(Q\) ラベル間遷移の数え上げ}（§3）: 3 色から 3 色への遷移は
  \(3 \times 3 = 9\) 通り。
\item
  \textbf{スピン 1 集合の数え上げ}（§4）: \(t\) 固定 2 状態、\(R\) 固定
  2 状態、\(Q\) 固定 9 状態、合計 13 状態。
\item
  \textbf{スピン 2 集合の 3 型}（§5）: \(tR, tQ, RQ\) の 3 通り。
\item
  \textbf{スピン 3 集合の 1 型}（§6）: \(tRQ\) の 1 通り。
\item
  \textbf{符号積 \(P(n) = (-1)^n\)}（定義 7.1、命題 7.1）:
  非零軸の符号の積は \(n\) の偶奇のみに依存する。
\item
  \textbf{符号付き面積 \(M(\sigma)\)}（定義 8.1）: \(R\) 軸を含む 32
  枚の正方形面から定義される本モデルの基本的な保存量。\(M(\sigma) = 0\)
  となる集合が存在する。
\end{enumerate}

これらはすべて本モデルの外部仮定（5 次元超直方体の採用、\(Q\)
軸の追加）の下で、組合せ論と幾何学から純粋に導出される結果である。

\begin{center}\rule{0.5\linewidth}{0.5pt}\end{center}

\subsection{10.
標準模型との対応の解釈例}\label{ux6a19ux6e96ux6a21ux578bux3068ux306eux5bfeux5fdcux306eux89e3ux91c8ux4f8b}

\textbf{本節の内容は本論文の主張ではなく、解釈の一例である。本モデルから標準模型を物理的に導出することを試みるものではない。}
本節では、§1--§9 で定義・導出した 6
次元構造とその数学的性質が、標準模型の粒子分類とどのように対応付けうるかの一例を示す。

以下の対応は本モデル内部からは正当化されず、\textbf{外部からの解釈の選択}である。将来、別の対応付けが可能になった場合、本節と並列に追加可能である。

\subsubsection{10.1
物理的ラベルの割り当て}\label{ux7269ux7406ux7684ux30e9ux30d9ux30ebux306eux5272ux308aux5f53ux3066}

§1.3.2 で述べた通り、5
軸は組合せ構造上すべて対等であるが、標準模型との対応付けにおいては以下のラベルを割り当てる。

\begin{itemize}
\tightlist
\item
  \(x_1, x_2, x_3 \to xyz\): 3 次元空間軸
\item
  \(x_4 \to t\): 時間軸
\item
  \(x_5 \to R\): 曲率軸（空間の曲率半径に相当）
\item
  \(Q\): 色ラベル軸（\(0\) = 色中性、\(r, g, b\) = 色荷）
\end{itemize}

\subsubsection{10.2
クォーク/レプトンへの対応}\label{ux30afux30a9ux30fcux30afux30ecux30d7ux30c8ux30f3ux3078ux306eux5bfeux5fdc}

\(Q\) 軸の値によって粒子の分類が以下のように対応する。

{\def\LTcaptype{none} % do not increment counter
\begin{longtable}[]{@{}cl@{}}
\toprule\noalign{}
\(Q\) 値 & 対応 \\
\midrule\noalign{}
\endhead
\bottomrule\noalign{}
\endlastfoot
\(0\) & レプトン（色荷なし） \\
\(r, g, b\) & クォーク（色荷あり） \\
\end{longtable}
}

同一の幾何学的集合 \((x, y, z, t, R)\) に対して、\(Q\) の値が 4
通り存在することにより、各世代のフェルミオンはレプトン 1 種とクォーク 3
色に分岐する。

\textbf{電荷の対応}:
電荷の符号は空間軸の集合符号によって区別される。クォークの分数電荷（\(\pm 1/3, \pm 2/3\)）の分母
3 は、\(Q\)
軸の非零値の数（\(|\{r, g, b\}| = 3\)）と一致する。電荷の完全な導出、特に
\(+2/3\) と \(-1/3\) の非対称性は今後の課題である（§10.10）。

\textbf{第 7 軸の検討}: 電荷を独立な第 7
軸として追加する設計も検討したが、非空間軸が 4 本に増えると \(n = 2\)
の型が \(\binom{4}{2} = 6\)
種に増加し、超対称性理論と類似の粒子スペクトルを生むことになる。現在の実験的状況（超対称性粒子の未発見）を踏まえ、本稿では電荷を
\(Q\) 軸の構造と空間軸の符号から読み取る方針を採る。

\subsubsection{10.3 グルーオン 8
状態への対応}\label{ux30b0ux30ebux30fcux30aaux30f3-8-ux72b6ux614bux3078ux306eux5bfeux5fdc}

§3 で数え上げた \(Q\) ラベル間遷移の 9
状態のうち、色シングレット（\(r\bar{r} + g\bar{g} + b\bar{b}\)）の 1
状態を除いた \textbf{8 状態をグルーオン \(g_1, \ldots, g_8\)}
に対応させる。

標準模型では、グルーオンの 8 状態は \(\mathrm{SU}(3)\) のリー代数
\(\mathfrak{su}(3)\) の 8
次元随伴表現として導出される。本解釈では、同じ数 8 が \(Q\) 軸の 4
ラベル構造から組合せ論的に出るが、\(\mathrm{SU}(3)\) の非可換構造定数
\(f^{abc}\) は本モデルからは導出されない。

\subsubsection{10.4 ゲージボソン 12
状態への対応}\label{ux30b2ux30fcux30b8ux30dcux30bdux30f3-12-ux72b6ux614bux3078ux306eux5bfeux5fdc}

§4 で数え上げたスピン 1 集合 13 状態のうち、\(Q\) 固定 9
状態から色シングレット 1 状態を除いた 8 状態と、\(t\) 固定 2 状態、\(R\)
固定 2 状態を合わせて \textbf{12 状態をゲージボソン} に対応させる。

{\def\LTcaptype{none} % do not increment counter
\begin{longtable}[]{@{}llc@{}}
\toprule\noalign{}
固定軸 & 対応 & 状態数 \\
\midrule\noalign{}
\endhead
\bottomrule\noalign{}
\endlastfoot
\(t\) & \(W^+, W^-\) & 2 \\
\(R\) & \(\gamma, Z^0\) & 2 \\
\(Q\) & \(g_1, \ldots, g_8\) & 8 \\
\textbf{合計} & & \textbf{12} \\
\end{longtable}
}

\(W\) ボソンの 2 状態（\(t = +1, -1\)）は荷電の正負に、\(\gamma\) と
\(Z^0\) は \(R\) 軸の符号による \(M(\sigma)\)
の有無（\(M(\gamma) = 0, M(Z) \neq 0\)）に対応する。合計 12
は、標準模型のゲージボソン状態数
\(\gamma + g \times 8 + W^+ + W^- + Z^0 = 12\) と一致する。

\subsubsection{10.5 スピン 2
集合の解釈}\label{ux30b9ux30d4ux30f3-2-ux96c6ux5408ux306eux89e3ux91c8}

§5 で数え上げたスピン 2 集合の 3 型に対して以下の対応を与える。

{\def\LTcaptype{none} % do not increment counter
\begin{longtable}[]{@{}
  >{\raggedright\arraybackslash}p{(\linewidth - 4\tabcolsep) * \real{0.2222}}
  >{\raggedright\arraybackslash}p{(\linewidth - 4\tabcolsep) * \real{0.4444}}
  >{\raggedright\arraybackslash}p{(\linewidth - 4\tabcolsep) * \real{0.3333}}@{}}
\toprule\noalign{}
\begin{minipage}[b]{\linewidth}\raggedright
型
\end{minipage} & \begin{minipage}[b]{\linewidth}\raggedright
固定軸
\end{minipage} & \begin{minipage}[b]{\linewidth}\raggedright
対応
\end{minipage} \\
\midrule\noalign{}
\endhead
\bottomrule\noalign{}
\endlastfoot
\(tR\) 型 & \(t, R\) & \textbf{グラビトン}（標準模型 +
重力の枠組みにおける対応） \\
\(tQ\) 型 & \(t, Q\) & 標準模型に対応する既知粒子なし \\
\(RQ\) 型 & \(R, Q\) & 標準模型に対応する既知粒子なし \\
\end{longtable}
}

命題 7.1 により、3 型すべてが \(P(2) = +1\) であり、§10.7
の解釈では引力のみを媒介する。\(tR\)
型については重力の性質（引力のみ）と整合する。\(tQ\) 型と \(RQ\)
型は、本モデルの組合せ構造から存在が導出される状態であるが、現在の標準模型には対応物がない。

\subsubsection{10.6 スピン 3
集合の解釈}\label{ux30b9ux30d4ux30f3-3-ux96c6ux5408ux306eux89e3ux91c8}

§6 で数え上げたスピン 3 集合の 1
型（\(tRQ\)）は、標準模型に対応する既知粒子が存在しない。

\subsubsection{\texorpdfstring{10.7 符号積 \(P(n)\)
と力の方向性の対応}{10.7 符号積 P(n) と力の方向性の対応}}\label{ux7b26ux53f7ux7a4d-pn-ux3068ux529bux306eux65b9ux5411ux6027ux306eux5bfeux5fdc}

§7 で定義した符号積 \(P(n) = (-1)^n\)
に対して、以下の物理的解釈を与える。

\begin{itemize}
\tightlist
\item
  \(P(n) = +1\)（\(n\) が偶数）: 力は一方向的（\textbf{引力のみ}）
\item
  \(P(n) = -1\)（\(n\) が奇数）:
  力は双方向的（\textbf{引力・斥力の両方}）
\end{itemize}

{\def\LTcaptype{none} % do not increment counter
\begin{longtable}[]{@{}ccc@{}}
\toprule\noalign{}
\(n\) & \(P(n)\) & 力の方向性の解釈 \\
\midrule\noalign{}
\endhead
\bottomrule\noalign{}
\endlastfoot
0 & \(+1\) & 方向なし（スカラー場、\(s = 0\)） \\
1 & \(-1\) & 引力・斥力両方（\(s = 1\)） \\
2 & \(+1\) & 引力のみ（\(s = 2\)） \\
3 & \(-1\) & 引力・斥力両方（\(s = 3\)） \\
\end{longtable}
}

この解釈の下で、スピン 2 集合の 3 型（\(tR, tQ, RQ\)）はすべて
\(n = 2\)（偶数）であり引力のみを媒介する。これは重力の性質（引力のみ）と整合する。スピン
1 の電磁力（\(n = 1\)、引力・斥力両方）とも整合する。

\textbf{注記}. この対応は本モデル内部からは正当化されず、\(P(n)\)
の偶奇と力の方向性の間の物理的機構は本稿では導出されない。

\subsubsection{\texorpdfstring{10.8 \(R\)
軸の値域と力の結合強度}{10.8 R 軸の値域と力の結合強度}}\label{r-ux8ef8ux306eux5024ux57dfux3068ux529bux306eux7d50ux5408ux5f37ux5ea6}

\textbf{\(R\) 軸の値域}: 先行研究 {[}W5{]}
において、離散空間での中心投影の整数論的条件から
\(R = \pm n^2 \cdot L_0\)（\(n \in \mathbb{N}\)、\(L_0\)
は単位長）の離散化が導かれている。ただし本稿の結果はこの離散化を必要とせず、\(R\)
が連続値をとる場合にも成立する。\(R\)
軸の値域を非有界と解釈すると、質量スペクトルの非有界性に対応する。

\textbf{各力と本モデルの軸の対応の解釈}:

{\def\LTcaptype{none} % do not increment counter
\begin{longtable}[]{@{}
  >{\raggedright\arraybackslash}p{(\linewidth - 6\tabcolsep) * \real{0.1053}}
  >{\raggedright\arraybackslash}p{(\linewidth - 6\tabcolsep) * \real{0.2632}}
  >{\raggedright\arraybackslash}p{(\linewidth - 6\tabcolsep) * \real{0.3158}}
  >{\raggedright\arraybackslash}p{(\linewidth - 6\tabcolsep) * \real{0.3158}}@{}}
\toprule\noalign{}
\begin{minipage}[b]{\linewidth}\raggedright
力
\end{minipage} & \begin{minipage}[b]{\linewidth}\raggedright
媒介粒子
\end{minipage} & \begin{minipage}[b]{\linewidth}\raggedright
関与する軸
\end{minipage} & \begin{minipage}[b]{\linewidth}\raggedright
結合の特徴
\end{minipage} \\
\midrule\noalign{}
\endhead
\bottomrule\noalign{}
\endlastfoot
強い力 & \(g_1\)--\(g_8\) & \(Q\) 軸（色遷移 8 チャンネル） &
色荷を持つ粒子間のみ \\
電磁力 & \(\gamma\) & \(R\) 軸 & 非零電荷を持つ粒子間 \\
弱い力 & \(W^\pm, Z^0\) & \(t\) 軸 & 全フェルミオン \\
重力 & \(G\) & \(t \wedge R\)（\(n = 2\)） & 全粒子（\(Q\) 非依存） \\
\end{longtable}
}

\textbf{注記}. 標準模型における結合定数は \(M_Z\) 付近で
\(\alpha_s \approx 0.1\), \(\alpha \approx 1/137\),
\(\alpha_W \approx 1/30\), \(\alpha_G \approx 10^{-39}\)
であり、強・電磁・弱は同程度のオーダー、重力のみ極端に弱い。この階層構造の定量的導出は今後の課題である。

\subsubsection{\texorpdfstring{10.9 \(M(\sigma)\)
と質量の類似性}{10.9 M(\textbackslash sigma) と質量の類似性}}\label{msigma-ux3068ux8ceaux91cfux306eux985eux4f3cux6027}

§8.1 で定義した符号付き面積 \(M(\sigma)\)
は、以下の点で粒子の質量と\textbf{類似性}を持つ。

\begin{itemize}
\tightlist
\item
  \(M(\sigma) = 0\)
  の集合は質量ゼロ粒子（光子、グルーオン）に対応する解釈が可能
\item
  \(R\) 軸の値域が非有界であることは、質量スペクトルの非有界性に対応する
\end{itemize}

ただし、本稿では \(M(\sigma)\)
と質量を同一視することはしない。具体的な質量値の導出は今後の課題である（§10.12）。

\subsubsection{10.10 表 A: 確認済み粒子と 6
次元集合の対応}\label{ux8868-a-ux78baux8a8dux6e08ux307fux7c92ux5b50ux3068-6-ux6b21ux5143ux96c6ux5408ux306eux5bfeux5fdc}

\textbf{ボソン}:

{\def\LTcaptype{none} % do not increment counter
\begin{longtable}[]{@{}
  >{\raggedright\arraybackslash}p{(\linewidth - 10\tabcolsep) * \real{0.1579}}
  >{\raggedright\arraybackslash}p{(\linewidth - 10\tabcolsep) * \real{0.1579}}
  >{\centering\arraybackslash}p{(\linewidth - 10\tabcolsep) * \real{0.1316}}
  >{\raggedright\arraybackslash}p{(\linewidth - 10\tabcolsep) * \real{0.1053}}
  >{\centering\arraybackslash}p{(\linewidth - 10\tabcolsep) * \real{0.1316}}
  >{\centering\arraybackslash}p{(\linewidth - 10\tabcolsep) * \real{0.3158}}@{}}
\toprule\noalign{}
\begin{minipage}[b]{\linewidth}\raggedright
粒子
\end{minipage} & \begin{minipage}[b]{\linewidth}\raggedright
記号
\end{minipage} & \begin{minipage}[b]{\linewidth}\centering
\(s\)
\end{minipage} & \begin{minipage}[b]{\linewidth}\raggedright
集合 \((x, y, z, t, R, Q)\)
\end{minipage} & \begin{minipage}[b]{\linewidth}\centering
\(n\)
\end{minipage} & \begin{minipage}[b]{\linewidth}\centering
質量 (MeV)
\end{minipage} \\
\midrule\noalign{}
\endhead
\bottomrule\noalign{}
\endlastfoot
ヒッグス & \(H^0\) & 0 & \((\pm, 0, 0, 0, 0, 0)\) 等 & 0 & 125,250 \\
\(W\) ボソン & \(W^+\) & 1 & \((0, 0, 0, +, 0, 0)\) & 1 & 80,379 \\
\(W\) ボソン & \(W^-\) & 1 & \((0, 0, 0, -, 0, 0)\) & 1 & 80,379 \\
光子 & \(\gamma\) & 1 & \((0, 0, 0, 0, +, 0)\) & 1 & 0 \\
\(Z\) ボソン & \(Z^0\) & 1 & \((0, 0, 0, 0, -, 0)\) & 1 & 91,188 \\
グルーオン & \(g_1\)--\(g_8\) & 1 & \((0, 0, 0, 0, 0, Q_i \to Q_j)\) & 1
& 0 \\
グラビトン & \(G\) & 2 & \((0, 0, 0, \pm, \pm, 0)\) & 2 & 0 \\
\end{longtable}
}

\textbf{レプトン} (\(Q = 0\)):

{\def\LTcaptype{none} % do not increment counter
\begin{longtable}[]{@{}
  >{\raggedright\arraybackslash}p{(\linewidth - 10\tabcolsep) * \real{0.1538}}
  >{\raggedright\arraybackslash}p{(\linewidth - 10\tabcolsep) * \real{0.1538}}
  >{\centering\arraybackslash}p{(\linewidth - 10\tabcolsep) * \real{0.1282}}
  >{\centering\arraybackslash}p{(\linewidth - 10\tabcolsep) * \real{0.1538}}
  >{\raggedright\arraybackslash}p{(\linewidth - 10\tabcolsep) * \real{0.1026}}
  >{\raggedright\arraybackslash}p{(\linewidth - 10\tabcolsep) * \real{0.3077}}@{}}
\toprule\noalign{}
\begin{minipage}[b]{\linewidth}\raggedright
粒子
\end{minipage} & \begin{minipage}[b]{\linewidth}\raggedright
記号
\end{minipage} & \begin{minipage}[b]{\linewidth}\centering
\(s\)
\end{minipage} & \begin{minipage}[b]{\linewidth}\centering
世代
\end{minipage} & \begin{minipage}[b]{\linewidth}\raggedright
集合
\end{minipage} & \begin{minipage}[b]{\linewidth}\raggedright
質量 (MeV)
\end{minipage} \\
\midrule\noalign{}
\endhead
\bottomrule\noalign{}
\endlastfoot
電子 & \(e^-\) & 1/2 & 1 & \((+, +, 0, 0, 0, 0)\) & 0.511 \\
電子ニュートリノ & \(\nu_e\) & 1/2 & 1 & \((+, -, 0, 0, 0, 0)\) &
\(< 10^{-6}\) \\
ミューオン & \(\mu^-\) & 1/2 & 2 & \((+, 0, +, 0, 0, 0)\) & 105.66 \\
ミューニュートリノ & \(\nu_\mu\) & 1/2 & 2 & \((+, 0, -, 0, 0, 0)\) &
\(< 0.17\) \\
タウ & \(\tau^-\) & 1/2 & 3 & \((0, +, +, 0, 0, 0)\) & 1,776.9 \\
タウニュートリノ & \(\nu_\tau\) & 1/2 & 3 & \((0, +, -, 0, 0, 0)\) &
\(< 15.5\) \\
\end{longtable}
}

反レプトンは空間軸の符号を反転した集合に対応する。

\textbf{クォーク} (\(Q \in \{r, g, b\}\)):

{\def\LTcaptype{none} % do not increment counter
\begin{longtable}[]{@{}llccll@{}}
\toprule\noalign{}
粒子 & 記号 & \(s\) & 世代 & 集合 & 質量 (MeV) \\
\midrule\noalign{}
\endhead
\bottomrule\noalign{}
\endlastfoot
アップ & \(u\) & 1/2 & 1 & \((+, +, 0, 0, 0, r/g/b)\) & 2.16 \\
ダウン & \(d\) & 1/2 & 1 & \((+, -, 0, 0, 0, r/g/b)\) & 4.67 \\
チャーム & \(c\) & 1/2 & 2 & \((+, 0, +, 0, 0, r/g/b)\) & 1,270 \\
ストレンジ & \(s\) & 1/2 & 2 & \((+, 0, -, 0, 0, r/g/b)\) & 93.4 \\
トップ & \(t\) & 1/2 & 3 & \((0, +, +, 0, 0, r/g/b)\) & 172,760 \\
ボトム & \(b\) & 1/2 & 3 & \((0, +, -, 0, 0, r/g/b)\) & 4,180 \\
\end{longtable}
}

各クォークは \(Q = r, g, b\) の 3
色状態を持つ。反クォークは空間軸の符号反転かつ \(Q\) の反色に対応する。

\textbf{集計（\(s = 0, 1/2, 1, 2\) の SM 対応可能な状態）}:

{\def\LTcaptype{none} % do not increment counter
\begin{longtable}[]{@{}lcccc@{}}
\toprule\noalign{}
分類 & \(s\) & \(Q\) & 本モデル導出数 & SM 対応 \\
\midrule\noalign{}
\endhead
\bottomrule\noalign{}
\endlastfoot
ヒッグス & 0 & 0 & 1 & 1 \\
ゲージボソン (\(t\) 型) & 1 & 0 & 2 & \(W^+, W^-\) \\
ゲージボソン (\(R\) 型) & 1 & 0 & 2 & \(\gamma, Z^0\) \\
グルーオン (\(Q\) 型) & 1 & 遷移 & 9 & 8（色シングレット 1 を除く） \\
グラビトン (\(tR\) 型) & 2 & 0 & 1 & \(G\)（SM 外） \\
レプトン & 1/2 & 0 & 6 & 6 \\
反レプトン & 1/2 & 0 & 6 & 6 \\
クォーク & 1/2 & \(r, g, b\) & 18 & 18 \\
反クォーク & 1/2 & \(\bar{r}, \bar{g}, \bar{b}\) & 18 & 18 \\
\textbf{小計} & & & \textbf{63} & \textbf{62（61 + G）} \\
\end{longtable}
}

（本モデル 63 状態 = SM 対応 62 状態 + 色シングレット 1 状態）。62
状態は標準模型 61 状態 + グラビトン 1 状態にマッピングされる。残り 1
状態（\(Q\)
遷移の色シングレット）は本モデルでは導出されるが、標準模型では独立なゲージボソンに対応しない。

上記に加え、\(s = 3/2, 2, 5/2, 3, 7/2\) の組合せ（表
B、§10.11）が本モデルから導出されるが、標準模型に対応する既知粒子は存在しない。

\subsubsection{10.11 表 B:
標準模型に対応物のない組合せ}\label{ux8868-b-ux6a19ux6e96ux6a21ux578bux306bux5bfeux5fdcux7269ux306eux306aux3044ux7d44ux5408ux305b}

本モデルの組合せ構造から導出されるが、標準模型に対応する既知粒子が存在しない状態を以下に列挙する。これらは本モデル内の組合せ論的帰結であり、未発見粒子の存在を主張するものではない。

\textbf{ボソン（\(xyz\) 全軸 \(0\)）}:

{\def\LTcaptype{none} % do not increment counter
\begin{longtable}[]{@{}
  >{\raggedright\arraybackslash}p{(\linewidth - 8\tabcolsep) * \real{0.1250}}
  >{\centering\arraybackslash}p{(\linewidth - 8\tabcolsep) * \real{0.1562}}
  >{\centering\arraybackslash}p{(\linewidth - 8\tabcolsep) * \real{0.1562}}
  >{\raggedright\arraybackslash}p{(\linewidth - 8\tabcolsep) * \real{0.3125}}
  >{\centering\arraybackslash}p{(\linewidth - 8\tabcolsep) * \real{0.2500}}@{}}
\toprule\noalign{}
\begin{minipage}[b]{\linewidth}\raggedright
型
\end{minipage} & \begin{minipage}[b]{\linewidth}\centering
\(s\)
\end{minipage} & \begin{minipage}[b]{\linewidth}\centering
\(n\)
\end{minipage} & \begin{minipage}[b]{\linewidth}\raggedright
集合形式
\end{minipage} & \begin{minipage}[b]{\linewidth}\centering
状態数
\end{minipage} \\
\midrule\noalign{}
\endhead
\bottomrule\noalign{}
\endlastfoot
\(tQ\) 型 & 2 & 2 & \((0, 0, 0, \pm, 0, Q_i \to Q_j)\) &
\(2 \times 9 = 18\) \\
\(RQ\) 型 & 2 & 2 & \((0, 0, 0, 0, \pm, Q_i \to Q_j)\) &
\(2 \times 9 = 18\) \\
\(tRQ\) 型 & 3 & 3 & \((0, 0, 0, \pm, \pm, Q_i \to Q_j)\) &
\(4 \times 9 = 36\) \\
\end{longtable}
}

\textbf{フェルミオン（\(xyz\) のうち 2 本非零、1 本 \(0\)）}:

{\def\LTcaptype{none} % do not increment counter
\begin{longtable}[]{@{}lcclcc@{}}
\toprule\noalign{}
型 & \(s\) & \(n\) & 非零の非空間軸 & 空間側パターン数 & \(Q\) 値 \\
\midrule\noalign{}
\endhead
\bottomrule\noalign{}
\endlastfoot
\(t\) 型 & 3/2 & 1 & \(t\) & \(3 \times 4 = 12\) & 4 \\
\(R\) 型 & 3/2 & 1 & \(R\) & 12 & 4 \\
\(Q\) 型 & 3/2 & 1 & \(Q\)（非零） & 12 & 3 \\
\(tR\) 型 & 5/2 & 2 & \(t, R\) & 12 & 4 \\
\(tQ\) 型 & 5/2 & 2 & \(t, Q\) & 12 & 3 \\
\(RQ\) 型 & 5/2 & 2 & \(R, Q\) & 12 & 3 \\
\(tRQ\) 型 & 7/2 & 3 & \(t, R, Q\) & 12 & 3 \\
\end{longtable}
}

\textbf{注記}. 空間側パターン数 12 は、零にする軸の選び方
\(\binom{3}{1} = 3\) 通り × 非零 2 軸の符号 \(2^2 = 4\)
通りに由来する。\(Q\) 値は \(Q = 0\) を含む場合 4 通り、\(Q\)
軸が非零の場合は \(\{r, g, b\}\) の 3 通りである。

\subsubsection{10.12
解釈例における未解決問題}\label{ux89e3ux91c8ux4f8bux306bux304aux3051ux308bux672aux89e3ux6c7aux554fux984c}

以下は本節の解釈の範囲における未解決問題である。

\textbf{10.12.1 \(M(\sigma)\) の定量計算と質量比}

符号付き面積 \(M(\sigma)\)
を各集合構造について明示的に計算し、既知の粒子質量比と比較することは、本解釈の検証手段である。

\begin{itemize}
\tightlist
\item
  \(m_\mu / m_e \approx 206.8\): 世代間の質量比
\item
  \(m_W / m_Z \approx 0.881\): ワインバーグ角との関係
\item
  \(m_t / m_e \approx 3.4 \times 10^5\): 最大質量比
\end{itemize}

\textbf{10.12.2 ゲージ群の導出}

\(Q\) 軸の 4 ラベル構造と \(\mathrm{SU}(3)\)、\(t\) 軸の 2 値と
\(\mathrm{SU}(2)\)、\(R\) 軸の連続スケールと \(\mathrm{U}(1)\)
の間の対応が示唆されるが、離散構造から連続的ゲージ群への移行の厳密な機構は未証明である。特にグルーオンの個数
8 は再現できるが、非可換構造定数 \(f^{abc}\) の情報は導出されない。

\textbf{10.12.3 CKM 行列と PMNS 行列}

世代間の混合（CKM 行列、PMNS 行列）が、\(S_3\)
対称性の破れとしてどのように導出されるかは未解決である。

\textbf{10.12.4 ヒッグス機構}

ヒッグスボソンは
\(n = 0\)（空間軸のみ）に対応するが、ヒッグス場による自発的対称性の破れがこの枠組みでどう理解されるかは今後の課題である。

\textbf{10.12.5 電荷値の完全な導出}

空間軸の集合符号から電荷符号が、\(Q\)
軸セクターから電荷の絶対値の単位（1 vs 1/3）が読み取れるが、\(+2/3\) と
\(-1/3\) の非対称性は本解釈の最も目立つ未解決問題である。

\textbf{10.12.6 \(s = 2, 3\) ボソンの追加予言}

\(tQ\) 型、\(RQ\) 型のスピン 2 ボソン、および \(tRQ\) 型のスピン 3
ボソンが存在するならば、その実験的帰結を明らかにする必要がある。

\subsubsection{10.13
解釈例の結論}\label{ux89e3ux91c8ux4f8bux306eux7d50ux8ad6}

本節では、§1--§9 で定義・導出した 6
次元構造の数学的性質を、標準模型の粒子分類と対応付ける一例を示した。

\begin{enumerate}
\def\labelenumi{\arabic{enumi}.}
\tightlist
\item
  スピン種数
  8（\(s = 0, 1/2, 1, 3/2, 2, 5/2, 3, 7/2\)）のうち、\(s = 0, 1/2, 1, 2\)
  の 63 状態のうち 62 状態が標準模型の既知粒子（61 +
  グラビトン）にマッピング可能（残り 1 状態は色シングレット）
\item
  \(Q\) 軸の 4 値が色荷とレプトン/クォーク分岐に対応
\item
  \(Q\) ラベル間遷移 9 通りのうち 8
  状態がグルーオンに対応（色シングレット 1 状態は対応物なし）
\item
  スピン 1 集合の 13 状態のうち 12 状態が標準模型のゲージボソンに対応
\item
  スピン 2 集合の \(tR\) 型がグラビトンに対応。\(tQ\) 型と \(RQ\)
  型は本モデルから導出されるが標準模型に対応物がない
\item
  スピン 3 集合の \(tRQ\) 型も同様に対応物がない
\item
  \(s = 3/2, 5/2, 7/2\)
  のフェルミオン型も組合せ論的に存在するが既知粒子に対応しない（表 B）
\item
  \(M(\sigma)\) が粒子質量と類似性を持つ
\item
  力ごとの関与軸の種類が結合強度の階層の定性的起源として解釈可能
\end{enumerate}

本モデルは標準模型の既知粒子を\textbf{超過}する組合せを導出する。これらの「余剰状態」は本解釈における未解決問題であり、(a)
未発見粒子の予言、(b) 物理的に実現されない組合せ（選択則の存在）、(c)
本解釈の限界、のいずれかとして今後の研究で検討される。

具体的な質量値、電荷値、ゲージ群の非可換構造、世代混合行列の導出は今後の課題である（§10.12）。

\begin{center}\rule{0.5\linewidth}{0.5pt}\end{center}

\subsection{参考文献}\label{ux53c2ux8003ux6587ux732e}

{[}1{]} Finkelstein, D. \& Rubinstein, J. ``Connection between Spin,
Statistics, and Kinks.'' \emph{J. Math. Phys.} 9, 1762 (1968).

{[}2{]} Atiyah, M.F., Manton, N.S. \& Schroers, B.J. ``Geometric Models
of Matter.'' \emph{Proc. R. Soc. A} 468, 1252 (2012). arXiv: 1108.5151.

{[}3{]} Furey, C. ``Standard Model Physics from an Algebra?'' arXiv:
1611.09182 (2016).

{[}4{]} Kaluza, Th. ``Zum Unitätsproblem der Physik.''
\emph{Sitzungsber. Preuss. Akad. Wiss. Berlin} 966--972 (1921).

{[}5{]} Coxeter, H.S.M. \emph{Regular Polytopes}. Dover (1973).

{[}6{]} Kihara, N. ``Bivectorial Classification of Spin Derived from the
Orientation Structure of the 5-Dimensional Hypercube.'' DOI:
\href{https://doi.org/10.5281/zenodo.19643358}{10.5281/zenodo.19643358}
(2026).

{[}7{]} Kihara, N. ``Geometric Classification of Spin Derived from the
Orientation Structure of the 5-Dimensional Orthoplex.'' DOI:
\href{https://doi.org/10.5281/zenodo.19630972}{10.5281/zenodo.19630972}
(2026).

{[}W5{]} Kihara, N. ``Complete Spherical Coverage of Central Projection
in Discrete Space and the Number-Theoretic Necessity of the
5-Dimensional Background Space.'' DOI:
\href{https://doi.org/10.5281/zenodo.19624957}{10.5281/zenodo.19624957}
(2026).

\end{document}
